\documentclass[10pt,a4paper]{amsart}
\usepackage{amsmath,amssymb,amsthm,graphicx}
% For commutative diagrams we use Xy-pic
\usepackage[all]{xy}

% We use 2cell for 2-commutative diagrams.
\xyoption{2cell}
% Theorem environments.
%
\theoremstyle{plain}
\newtheorem{theorem}[subsection]{Theorem}
\newtheorem{proposition}[subsection]{Proposition}
\newtheorem{lemma}[subsection]{Lemma}

\theoremstyle{definition}
\newtheorem{definition}[subsection]{Definition}
\newtheorem{example}[subsection]{Example}
\newtheorem{exercise}[subsection]{Exercise}
\newtheorem{situation}[subsection]{Situation}

\theoremstyle{remark}
\newtheorem{remark}[subsection]{Remark}
\newtheorem{remarks}[subsection]{Remarks}

\numberwithin{equation}{subsection}

% Macros
%
\def\lim{\mathop{\mathrm{lim}}\nolimits}
\def\colim{\mathop{\mathrm{colim}}\nolimits}
\def\Spec{\mathop{\mathrm{Spec}}}
\def\Hom{\mathop{\mathrm{Hom}}\nolimits}
\def\Ext{\mathop{\mathrm{Ext}}\nolimits}
\def\SheafHom{\mathop{\mathcal{H}\!\mathit{om}}\nolimits}
\def\SheafExt{\mathop{\mathcal{E}\!\mathit{xt}}\nolimits}
\def\Sch{\mathit{Sch}}
\def\Mor{\mathop{\mathrm{Mor}}\nolimits}
\def\Ob{\mathop{\mathrm{Ob}}\nolimits}
\def\Sh{\mathop{\mathit{Sh}}\nolimits}
\def\NL{\mathop{N\!L}\nolimits}
\def\CH{\mathop{\mathrm{CH}}\nolimits}
\def\proetale{{pro\text{-}\acute{e}tale}}
\def\etale{{\acute{e}tale}}
\def\QCoh{\mathit{QCoh}}
\def\Ker{\mathop{\mathrm{Ker}}}
\def\Im{\mathop{\mathrm{Im}}}
\def\Coker{\mathop{\mathrm{Coker}}}
\def\Coim{\mathop{\mathrm{Coim}}}

% Boxtimes
%
\DeclareMathSymbol{\boxtimes}{\mathbin}{AMSa}{"02}

%
% Macros for moduli stacks/spaces
%
\def\QCohstack{\mathcal{QC}\!\mathit{oh}}
\def\Cohstack{\mathcal{C}\!\mathit{oh}}
\def\Spacesstack{\mathcal{S}\!\mathit{paces}}
\def\Quotfunctor{\mathrm{Quot}}
\def\Hilbfunctor{\mathrm{Hilb}}
\def\Curvesstack{\mathcal{C}\!\mathit{urves}}
\def\Polarizedstack{\mathcal{P}\!\mathit{olarized}}
\def\Complexesstack{\mathcal{C}\!\mathit{omplexes}}
% \Pic is the operator that assigns to X its picard group, usage \Pic(X)
% \Picardstack_{X/B} denotes the Picard stack of X over B
% \Picardfunctor_{X/B} denotes the Picard functor of X over B
\def\Pic{\mathop{\mathrm{Pic}}\nolimits}
\def\Picardstack{\mathcal{P}\!\mathit{ic}}
\def\Picardfunctor{\mathrm{Pic}}
\def\Deformationcategory{\mathcal{D}\!\mathit{ef}}
\usepackage{cleveref}
\crefname{originalreference}{original reference}{original references}
\Crefname{originalreference}{Original reference}{Original references}
\makeatletter
\newlabel{lemma-blow-up-square-cohomological-descent}{{0EW5}{source}{}{sourcepointer-1}{}}
\newlabel{lemma-blow-up-square-cohomological-descent@cref}{{[originalreference][0][]0EW5}{[1][0][]source}{}{}{}}
\makeatother
\crefname{originalreference}{original reference}{original references}
\Crefname{originalreference}{Original reference}{Original references}
\makeatletter
\newlabel{lemma-proper-base-change}{{09V5}{source}{}{sourcepointer-2}{}}
\newlabel{lemma-proper-base-change@cref}{{[originalreference][0][]09V5}{[1][0][]source}{}{}{}}
\makeatother
\title{Descent along the blow-up square: compatible data on the centre and on the blow-up glue to an object of the derived category of the ambient scheme}
\author{The Stacks Project Authors}
\date{stacks-project:etale-cohomology.tex:local-snapshot-20261011; GNU FDL 1.2 or later}
\begin{document}
\maketitle
\noindent\textit{Editorial note.} All delivered mathematical text is the author's own, copied from the pinned Stacks Project chapter file etale-cohomology.tex; no statement or proof step was rewritten, reordered or removed.\par\medskip
\noindent\textit{Editorial note.} The counted proof is the author's complete argument: the proof is written out in full: the construction of the distinguished triangle defining K, the definition of the two structure maps by adjunction and by gamma, the two isomorphism arguments for g and for f, and the final map-of-triangles argument for the morphism part.\par\medskip
\noindent\textit{Editorial note.} The label lemma-blow-up-square-equivalence names the counted statement, and the proof environment immediately following it in the source is delivered verbatim as its proof.\par\medskip
\noindent\textit{Editorial note.} The counted statement is carried by the author's own lemma environment in the chapter source.\par\medskip
\noindent\textit{Editorial note.} The statement and proof are copied from the pinned chapter file; the two cited inputs are kept as external pointers with their canonical Stacks tags.\par\medskip
\noindent\textit{Editorial note.} The preamble is the chapter's own native preamble excerpt already used for this source: the Xy-pic setup (native preamble lines 14-18) and the native theorem declarations with the mathematical macros (native preamble lines 154-217).\par\medskip
\noindent\textit{Allowed standard prerequisites (editorial).}\begin{itemize}
\item The blow-up square of a closed subscheme cut out by a quasi-coherent ideal of finite type. Same chapter, the blow-up square with E = Z x\_X X' and the equality c = i o pi = b o j of the composite morphisms. Supplies the geometric set-up and the four morphisms entering the constructed triangle.
\item Cohomological descent triangle of the blow-up square. Same chapter, the distinguished triangle K -> Ri\_*(K|\_Z) + Rb\_*(K|\_\{X'\}) -> Rc\_*(K|\_E) -> K[1] for torsion complexes that the constructed object must fit in. Supplies the shape of the triangle defining K, and the reason the construction is the only possible one.
\item Proper base change for the blow-up and adjunction for the pushforwards along i, b, j. Same chapter, the proper base change theorem and the adjunctions relating pullback and pushforward along the morphisms of the square. Supplies the identification making g and f isomorphisms after restricting the triangle to Z and to X'.
\end{itemize}\medskip
\ifdefined\hypertarget\hypertarget{sourcepointer-1}{}\fi
\noindent\textit{Source reference.} Original reference 0EW5: sources/stacks-etale-cohomology/pinned-original-source.tex, source lines 23315--23315. Prerequisite pointer invoked at proof line(s) 23472; it names a standard fact of the same chapter that the counted argument uses, and every step of the counted proof built on it is retained verbatim in the delivered body.\par\medskip
\ifdefined\hypertarget\hypertarget{sourcepointer-2}{}\fi
\noindent\textit{Source reference.} Original reference 09V5: sources/stacks-etale-cohomology/pinned-original-source.tex, source lines 18708--18708. Prerequisite pointer invoked at proof line(s) 23485; it names a standard fact of the same chapter that the counted argument uses, and every step of the counted proof built on it is retained verbatim in the delivered body.\par\medskip
\begin{lemma}
\label{lemma-blow-up-square-equivalence}
Let $X$ be a scheme and let $Z \subset X$ be a closed subscheme
cut out by a quasi-coherent ideal of finite type. Consider the
corresponding blow up square
$$
\xymatrix{
E \ar[d]_\pi \ar[r]_j & X' \ar[d]^b \\
Z \ar[r]^i & X
}
$$
Suppose given
\begin{enumerate}
\item an object $K'$ of $D^+(X'_\etale)$ with torsion cohomology sheaves,
\item an object $L$ of $D^+(Z_\etale)$ with torsion cohomology sheaves, and
\item an isomorphism $\gamma : K'|_E \to L|_E$.
\end{enumerate}
Then there exists an object $K$ of $D^+(X_\etale)$
and isomorphisms $f : K|_{X'} \to K'$, $g : K|_Z \to L$ such
that $\gamma = g|_E \circ f^{-1}|_E$.
Moreover, given
\begin{enumerate}
\item an object $M$ of $D^+(X_\etale)$ with torsion cohomology sheaves,
\item a morphism $\alpha : K' \to M|_{X'}$ of $D(X'_\etale)$,
\item a morphism $\beta : L \to M|_Z$ of $D(Z_\etale)$,
\end{enumerate}
such that
$$
\alpha|_E  = \beta|_E \circ \gamma.
$$
Then there exists a morphism $M \to K$ in $D(X_\etale)$
whose restriction to $X'$ is $a \circ f$
and whose restriction to $Z$ is $b \circ g$.
\end{lemma}

\begin{proof}
If $K$ exists, then Lemma \ref{lemma-blow-up-square-cohomological-descent}
tells us a distinguished triangle that it fits in. Thus we simply choose
a distinguished triangle
$$
K \to Ri_*(L) \oplus Rb_*(K') \to Rc_*(L|_E) \to K[1]
$$
where $c = i \circ \pi = b \circ j$. Here the map $Ri_*(L) \to Rc_*(L|_E)$
is $Ri_*$ applied to the adjunction mapping $E \to R\pi_*(L|_E)$.
The map $Rb_*(K') \to Rc_*(L|_E)$ is the composition of the canonical map
$Rb_*(K') \to Rc_*(K'|_E)) = R$ and $Rc_*(\gamma)$.
The maps $g$ and $f$ of the statement of the lemma are the adjoints
of these maps. If we restrict this distinguished triangle to $Z$
then the map $Rb_*(K) \to Rc_*(L|_E)$ becomes an isomorphism
by the base change theorem (Lemma \ref{lemma-proper-base-change}) and hence
the map $g : K|_Z \to L$ is an isomorphism.
Looking at the distinguished triangle we see that $f : K|_{X'} \to K'$
is an isomorphism over $X' \setminus E = X \setminus Z$.
Moreover, we have $\gamma \circ f|_E = g|_E$ by construction.
Then since $\gamma$ and $g$ are isomorphisms we conclude
that $f$ induces isomorphisms on stalks at geometric points of $E$
as well. Thus $f$ is an isomorphism.

\medskip\noindent
For the final statement, we may replace $K'$ by $K|_{X'}$,
$L$ by $K|_Z$, and $\gamma$ by the canonical identification.
Observe that $\alpha$ and $\beta$ induce a commutative square
$$
\xymatrix{
K \ar[r] \ar@{..>}[d] &
Ri_*(K|_Z) \oplus Rb_*(K|_{X'}) \ar[r] \ar[d]_{\beta \oplus \alpha} &
Rc_*(K|_E) \ar[r] \ar[d]_{\alpha|_E} &
K[1] \ar@{..>}[d] \\
M \ar[r] &
Ri_*(M|_Z) \oplus Rb_*(M|_{X'}) \ar[r] &
Rc_*(M|_E) \ar[r] &
M[1]
}
$$
Thus by the axioms of a derived category we get a dotted
arrow producing a morphism of distinguished triangles.
\end{proof}

\end{document}
